\chapter{Visualization Dashboard}

\section{Overview}

The visualization dashboard provides 10 executive-level visualizations for payroll analysis:

\begin{enumerate}
    \item Total Cost Trend - Monthly spending over time
    \item Cost Forecast - 6-month prediction using AR(2)
    \item Cost Per Employee - Efficiency metric
    \item Cost Breakdown - Salary component distribution
    \item Cost Distribution - Box plot \& histogram for equity
    \item Top Employees - 10 most expensive (retention risks)
    \item Gross vs Net Comparison - Interactive employee selection
    \item Headcount Trend - Workforce growth tracking
    \item Burn Rate Dashboard - 4-panel cash management
    \item Year-over-Year Comparison - Seasonal analysis
\end{enumerate}

\section{Architecture}

\subsection{Class Structure}

\begin{lstlisting}[style=pythonstyle]
class PayrollDashboard:
    """Interactive tkinter GUI for payroll visualizations."""
    
    def __init__(self, data_path):
        self.df = pd.read_csv(data_path)
        self.root = tk.Tk()
        self.root.title("Payroll Analytics Dashboard")
        
        # Create GUI
        self._create_title_label()
        self._create_plot_buttons()
        self._configure_layout()
    
    def _create_plot_buttons(self):
        """Create 10 buttons for different plots."""
        buttons = [
            ("1. Total Cost Trend", self.plot_total_cost_trend),
            ("2. Cost Forecast (AR)", self.plot_cost_forecast),
            ("3. Cost Per Employee", self.plot_cost_per_employee),
            # ... 7 more buttons ...
        ]
        
        for text, command in buttons:
            btn = tk.Button(
                self.root, text=text, command=command,
                width=30, height=2, font=('Arial', 11)
            )
            btn.pack(pady=5)
    
    def run(self):
        """Start the GUI event loop."""
        self.root.mainloop()
\end{lstlisting}

\subsection{Data Requirements}

\textbf{Input Format:} Long-format CSV with columns:
\begin{itemize}
    \item \texttt{month} - Date (YYYY-MM-DD)
    \item \texttt{employee\_id} - Unique identifier
    \item \texttt{salaire\_brut} - Gross salary
    \item \texttt{salaire\_net} - Net salary
    \item Additional salary components (optional)
\end{itemize}

\textbf{Preprocessing:}
\begin{lstlisting}[style=pythonstyle]
def __init__(self, data_path):
    # Load data
    self.df = pd.read_csv(data_path)
    
    # Convert month to datetime
    self.df['month'] = pd.to_datetime(self.df['month'])
    
    # Sort by month
    self.df = self.df.sort_values('month')
\end{lstlisting}

\section{Plot 1: Total Cost Trend}

\subsection{Purpose}
Monitor monthly payroll spending to track budget usage and identify unusual spikes.

\subsection{Audience}
\begin{itemize}
    \item \textbf{CEO:} Overall spending trends
    \item \textbf{CFO:} Budget compliance
    \item \textbf{HR:} Cost impact of hiring/terminations
\end{itemize}

\subsection{Implementation}

\begin{lstlisting}[style=pythonstyle]
def plot_total_cost_trend(self):
    """Line plot of total monthly payroll cost."""
    
    # Aggregate: sum salaire_brut per month
    monthly = self.df.groupby('month')['salaire_brut'].sum()
    
    # Create figure
    fig, ax = plt.subplots(figsize=(12, 6))
    
    # Plot line with markers
    ax.plot(monthly.index, monthly.values, 
            marker='o', linewidth=2, markersize=8,
            color='#2E86AB', label='Total Cost')
    
    # Formatting
    ax.set_title('Total Monthly Payroll Cost', 
                 fontsize=16, weight='bold')
    ax.set_xlabel('Month', fontsize=12)
    ax.set_ylabel('Total Cost (MAD)', fontsize=12)
    ax.grid(True, alpha=0.3)
    ax.legend()
    
    # Format y-axis as currency
    ax.yaxis.set_major_formatter(
        plt.FuncFormatter(lambda x, p: f'{x:,.0f} MAD')
    )
    
    # Rotate x-axis labels
    plt.xticks(rotation=45, ha='right')
    plt.tight_layout()
    plt.show()
\end{lstlisting}

\subsection{Insights}

\begin{itemize}
    \item \textbf{Upward trend:} Hiring or salary increases
    \item \textbf{Sudden spike:} Bonuses or back pay
    \item \textbf{Downward trend:} Layoffs or cost reduction
    \item \textbf{Seasonal pattern:} Year-end bonuses
\end{itemize}

\section{Plot 2: Cost Forecast (AR Model)}

\subsection{Purpose}
Predict next 6 months of payroll spending using autoregressive model for budget planning.

\subsection{Methodology: AR(2) Model}

\textbf{Theory:}
$$X_t = c + \phi_1 X_{t-1} + \phi_2 X_{t-2} + \epsilon_t$$

Where:
\begin{itemize}
    \item $X_t$ = Total cost at time $t$
    \item $\phi_1, \phi_2$ = Autoregressive coefficients
    \item $c$ = Constant term
    \item $\epsilon_t$ = White noise error
\end{itemize}

\textbf{Why AR(2)?} User validated through testing that AR(2) gives most accurate predictions for this payroll dataset.

\subsection{Implementation}

\begin{lstlisting}[style=pythonstyle]
def plot_cost_forecast(self):
    """Forecast next 6 months using AR(2) model."""
    from statsmodels.tsa.ar_model import AutoReg
    
    # Prepare time series
    monthly = self.df.groupby('month')['salaire_brut'].sum()
    
    # Fit AR(2) model
    model = AutoReg(monthly, lags=2)
    model_fit = model.fit()
    
    # Forecast 6 steps ahead
    forecast_steps = 6
    predictions = model_fit.forecast(steps=forecast_steps)
    
    # Generate future dates
    last_date = monthly.index[-1]
    future_dates = pd.date_range(
        start=last_date + pd.DateOffset(months=1),
        periods=forecast_steps,
        freq='MS'
    )
    
    # Calculate confidence intervals (95%)
    # Variance grows with horizon: sigma^2 * (1 + phi1^2 + phi2^2)
    std_error = model_fit.params.iloc[-1]  # Residual std
    ci_lower = predictions - 1.96 * std_error
    ci_upper = predictions + 1.96 * std_error
    
    # Plot
    fig, ax = plt.subplots(figsize=(14, 6))
    
    # Historical data
    ax.plot(monthly.index, monthly.values,
            marker='o', linewidth=2, label='Historical',
            color='#2E86AB')
    
    # Forecast
    ax.plot(future_dates, predictions,
            marker='s', linewidth=2, linestyle='--',
            label='Forecast (AR2)', color='#A23B72')
    
    # Confidence interval
    ax.fill_between(future_dates, ci_lower, ci_upper,
                     alpha=0.2, color='#A23B72',
                     label='95% Confidence Interval')
    
    # Formatting
    ax.set_title('6-Month Cost Forecast (AR2 Model)',
                 fontsize=16, weight='bold')
    ax.set_xlabel('Month', fontsize=12)
    ax.set_ylabel('Total Cost (MAD)', fontsize=12)
    ax.grid(True, alpha=0.3)
    ax.legend()
    
    plt.xticks(rotation=45, ha='right')
    plt.tight_layout()
    plt.show()
\end{lstlisting}

\subsection{Interpretation}

\begin{table}[H]
\centering
\begin{tabular}{ll}
\toprule
\textbf{Element} & \textbf{Meaning} \\
\midrule
Forecast line & Point estimate for each month \\
Shaded region & 95\% probability cost falls in range \\
Widening region & Uncertainty increases with time \\
\bottomrule
\end{tabular}
\caption{Forecast Plot Elements}
\end{table}

\section{Plot 3: Cost Per Employee}

\subsection{Purpose}
Track efficiency metric (average cost per employee) to identify staffing efficiency.

\subsection{Implementation}

\begin{lstlisting}[style=pythonstyle]
def plot_cost_per_employee(self):
    """Line plot of average cost per employee."""
    
    # Calculate monthly: total_cost / num_employees
    monthly_stats = self.df.groupby('month').agg({
        'salaire_brut': 'sum',
        'employee_id': 'nunique'  # Distinct count
    })
    
    monthly_stats['cost_per_emp'] = (
        monthly_stats['salaire_brut'] / 
        monthly_stats['employee_id']
    )
    
    # Plot
    fig, ax = plt.subplots(figsize=(12, 6))
    ax.plot(monthly_stats.index, 
            monthly_stats['cost_per_emp'],
            marker='o', linewidth=2, color='#F18F01')
    
    # Add mean line
    mean_cost = monthly_stats['cost_per_emp'].mean()
    ax.axhline(mean_cost, color='red', linestyle='--',
               label=f'Mean: {mean_cost:,.0f} MAD')
    
    ax.set_title('Average Cost Per Employee',
                 fontsize=16, weight='bold')
    ax.set_xlabel('Month', fontsize=12)
    ax.set_ylabel('Cost Per Employee (MAD)', fontsize=12)
    ax.grid(True, alpha=0.3)
    ax.legend()
    
    plt.xticks(rotation=45, ha='right')
    plt.tight_layout()
    plt.show()
\end{lstlisting}

\subsection{Insights}

\begin{itemize}
    \item \textbf{Increasing:} Hiring more senior/expensive staff
    \item \textbf{Decreasing:} Hiring junior staff or salary cuts
    \item \textbf{Stable:} Consistent workforce composition
\end{itemize}

\section{Plot 4: Cost Breakdown (Stacked Area)}

\subsection{Purpose}
Visualize composition of payroll (gross, net, taxes) over time.

\subsection{Implementation}

\begin{lstlisting}[style=pythonstyle]
def plot_cost_breakdown(self):
    """Stacked area chart of salary components."""
    
    # Aggregate multiple columns
    monthly = self.df.groupby('month').agg({
        'salaire_brut': 'sum',
        'salaire_net': 'sum',
        # Add other components if available
    }).fillna(0)
    
    # Calculate derived: taxes = brut - net
    monthly['taxes_total'] = (
        monthly['salaire_brut'] - monthly['salaire_net']
    )
    
    # Plot stacked area
    fig, ax = plt.subplots(figsize=(14, 7))
    ax.stackplot(monthly.index,
                 monthly['salaire_net'],
                 monthly['taxes_total'],
                 labels=['Net Salary', 'Taxes & Contributions'],
                 colors=['#06A77D', '#D62246'],
                 alpha=0.8)
    
    ax.set_title('Payroll Cost Breakdown',
                 fontsize=16, weight='bold')
    ax.set_xlabel('Month', fontsize=12)
    ax.set_ylabel('Amount (MAD)', fontsize=12)
    ax.legend(loc='upper left')
    ax.grid(True, alpha=0.3)
    
    plt.xticks(rotation=45, ha='right')
    plt.tight_layout()
    plt.show()
\end{lstlisting}

\section{Plot 5: Cost Distribution}

\subsection{Purpose}
Analyze salary equity using box plot (outliers) and histogram (distribution shape).

\subsection{Implementation}

\begin{lstlisting}[style=pythonstyle]
def plot_cost_distribution(self):
    """Box plot + histogram for salary distribution."""
    
    # Get latest month data
    latest_month = self.df['month'].max()
    latest_data = self.df[self.df['month'] == latest_month]
    
    # Create subplots
    fig, (ax1, ax2) = plt.subplots(1, 2, figsize=(14, 6))
    
    # Box plot
    ax1.boxplot(latest_data['salaire_brut'], vert=True)
    ax1.set_title('Salary Distribution (Box Plot)',
                  fontsize=14, weight='bold')
    ax1.set_ylabel('Gross Salary (MAD)', fontsize=12)
    ax1.grid(True, alpha=0.3, axis='y')
    
    # Histogram
    ax2.hist(latest_data['salaire_brut'], bins=15,
             color='#577590', alpha=0.7, edgecolor='black')
    ax2.set_title('Salary Distribution (Histogram)',
                  fontsize=14, weight='bold')
    ax2.set_xlabel('Gross Salary (MAD)', fontsize=12)
    ax2.set_ylabel('Number of Employees', fontsize=12)
    ax2.grid(True, alpha=0.3, axis='y')
    
    plt.tight_layout()
    plt.show()
\end{lstlisting}

\subsection{Insights}

\begin{description}
    \item[Box Plot] Shows quartiles and outliers
    \begin{itemize}
        \item Points beyond whiskers = outliers
        \item Wide box = high salary variance
    \end{itemize}
    
    \item[Histogram] Shows distribution shape
    \begin{itemize}
        \item Right-skewed = few high earners
        \item Normal = balanced distribution
    \end{itemize}
\end{description}

\section{Plot 6: Top 10 Employees}

\subsection{Purpose}
Identify retention risks (most expensive employees) for succession planning.

\subsection{Implementation}

\begin{lstlisting}[style=pythonstyle]
def plot_top_employees(self):
    """Horizontal bar chart of top 10 most expensive employees."""
    
    # Average salary per employee across all months
    emp_avg = self.df.groupby('employee_id')['salaire_brut'].mean()
    
    # Get top 10
    top10 = emp_avg.nlargest(10).sort_values()
    
    # Plot horizontal bar
    fig, ax = plt.subplots(figsize=(10, 8))
    bars = ax.barh(range(len(top10)), top10.values,
                   color='#C73E1D', alpha=0.8)
    
    # Customize
    ax.set_yticks(range(len(top10)))
    ax.set_yticklabels(top10.index)
    ax.set_xlabel('Average Gross Salary (MAD)', fontsize=12)
    ax.set_title('Top 10 Most Expensive Employees',
                 fontsize=16, weight='bold')
    ax.grid(True, alpha=0.3, axis='x')
    
    # Add value labels
    for i, bar in enumerate(bars):
        width = bar.get_width()
        ax.text(width, i, f' {width:,.0f}',
                va='center', fontsize=10)
    
    plt.tight_layout()
    plt.show()
\end{lstlisting}

\section{Plot 7: Gross vs Net (Interactive)}

\subsection{Purpose}
Compare gross and net salary for individual employee to analyze tax burden.

\subsection{Interactive Selection}

\begin{lstlisting}[style=pythonstyle]
def plot_gross_vs_net(self):
    """Interactive plot: user selects employee."""
    
    # Get unique employees
    employees = sorted(self.df['employee_id'].unique())
    
    # Create selection dialog
    selection_window = tk.Toplevel(self.root)
    selection_window.title("Select Employee")
    selection_window.geometry("300x400")
    
    tk.Label(selection_window, 
             text="Select an employee:",
             font=('Arial', 12, 'bold')).pack(pady=10)
    
    # Listbox with scrollbar
    frame = tk.Frame(selection_window)
    frame.pack(fill=tk.BOTH, expand=True, padx=10, pady=10)
    
    scrollbar = tk.Scrollbar(frame)
    scrollbar.pack(side=tk.RIGHT, fill=tk.Y)
    
    listbox = tk.Listbox(frame, yscrollcommand=scrollbar.set,
                        font=('Arial', 10))
    listbox.pack(side=tk.LEFT, fill=tk.BOTH, expand=True)
    scrollbar.config(command=listbox.yview)
    
    # Populate listbox
    for emp in employees:
        listbox.insert(tk.END, emp)
    
    # Plot button
    def on_plot():
        selection = listbox.curselection()
        if not selection:
            messagebox.showwarning("No Selection", 
                                  "Please select an employee")
            return
        
        emp_id = employees[selection[0]]
        self._plot_gross_vs_net_for_employee(emp_id)
        selection_window.destroy()
    
    tk.Button(selection_window, text="Plot",
             command=on_plot, width=15, height=2).pack(pady=10)
\end{lstlisting}

\subsection{Plotting Function}

\begin{lstlisting}[style=pythonstyle]
def _plot_gross_vs_net_for_employee(self, emp_id):
    """Plot gross vs net for specific employee."""
    
    # Filter data
    emp_data = self.df[self.df['employee_id'] == emp_id].sort_values('month')
    
    if emp_data.empty:
        messagebox.showerror("Error", f"No data for {emp_id}")
        return
    
    # Plot
    fig, ax = plt.subplots(figsize=(12, 6))
    
    x = emp_data['month']
    ax.plot(x, emp_data['salaire_brut'], 
            marker='o', label='Gross', color='#2E86AB', linewidth=2)
    ax.plot(x, emp_data['salaire_net'],
            marker='s', label='Net', color='#06A77D', linewidth=2)
    
    # Fill area between (tax burden)
    ax.fill_between(x, emp_data['salaire_brut'],
                     emp_data['salaire_net'],
                     alpha=0.3, color='#D62246',
                     label='Taxes & Contributions')
    
    ax.set_title(f'Gross vs Net Salary: {emp_id}',
                 fontsize=16, weight='bold')
    ax.set_xlabel('Month', fontsize=12)
    ax.set_ylabel('Salary (MAD)', fontsize=12)
    ax.legend()
    ax.grid(True, alpha=0.3)
    
    plt.xticks(rotation=45, ha='right')
    plt.tight_layout()
    plt.show()
\end{lstlisting}

\section{Plot 8: Headcount Trend}

\subsection{Implementation}

\begin{lstlisting}[style=pythonstyle]
def plot_headcount_trend(self):
    """Line plot of employee count over time."""
    
    # Count unique employees per month
    headcount = self.df.groupby('month')['employee_id'].nunique()
    
    fig, ax = plt.subplots(figsize=(12, 6))
    ax.plot(headcount.index, headcount.values,
            marker='o', linewidth=2, color='#06A77D')
    
    ax.set_title('Employee Headcount Trend',
                 fontsize=16, weight='bold')
    ax.set_xlabel('Month', fontsize=12)
    ax.set_ylabel('Number of Employees', fontsize=12)
    ax.grid(True, alpha=0.3)
    
    # Add annotations for significant changes
    changes = headcount.diff()
    for date, change in changes.items():
        if abs(change) >= 5:  # Significant change threshold
            label = f'+{int(change)}' if change > 0 else f'{int(change)}'
            ax.annotate(label, xy=(date, headcount[date]),
                       xytext=(0, 10), textcoords='offset points',
                       ha='center', fontsize=9, color='red',
                       bbox=dict(boxstyle='round,pad=0.3',
                                facecolor='yellow', alpha=0.7))
    
    plt.xticks(rotation=45, ha='right')
    plt.tight_layout()
    plt.show()
\end{lstlisting}

\section{Plot 9: Burn Rate Dashboard}

\subsection{Purpose}
4-panel dashboard for cash management (monthly cost, cumulative, moving average, runway).

\subsection{Implementation}

\begin{lstlisting}[style=pythonstyle]
def plot_burn_rate(self):
    """4-panel dashboard for cash management."""
    
    monthly = self.df.groupby('month')['salaire_brut'].sum()
    
    fig, ((ax1, ax2), (ax3, ax4)) = plt.subplots(2, 2, figsize=(16, 10))
    
    # Panel 1: Monthly cost
    ax1.bar(monthly.index, monthly.values, color='#577590', alpha=0.7)
    ax1.set_title('Monthly Payroll Cost', fontsize=14, weight='bold')
    ax1.set_ylabel('Cost (MAD)', fontsize=11)
    ax1.grid(True, alpha=0.3, axis='y')
    plt.setp(ax1.xaxis.get_majorticklabels(), rotation=45, ha='right')
    
    # Panel 2: Cumulative cost
    cumulative = monthly.cumsum()
    ax2.plot(cumulative.index, cumulative.values,
             marker='o', linewidth=2, color='#D62246')
    ax2.set_title('Cumulative Cost', fontsize=14, weight='bold')
    ax2.set_ylabel('Total Cost (MAD)', fontsize=11)
    ax2.grid(True, alpha=0.3)
    plt.setp(ax2.xaxis.get_majorticklabels(), rotation=45, ha='right')
    
    # Panel 3: 3-month moving average
    ma3 = monthly.rolling(window=3).mean()
    ax3.plot(monthly.index, monthly.values,
             alpha=0.3, label='Monthly', color='gray')
    ax3.plot(ma3.index, ma3.values,
             linewidth=2, label='3-Month MA', color='#F18F01')
    ax3.set_title('3-Month Moving Average', fontsize=14, weight='bold')
    ax3.set_ylabel('Cost (MAD)', fontsize=11)
    ax3.legend()
    ax3.grid(True, alpha=0.3)
    plt.setp(ax3.xaxis.get_majorticklabels(), rotation=45, ha='right')
    
    # Panel 4: Runway (assuming fixed budget)
    budget = 500000  # Example budget
    months_left = budget / monthly.mean()
    ax4.text(0.5, 0.5, f'Runway:\n{months_left:.1f} months',
             ha='center', va='center', fontsize=24, weight='bold')
    ax4.text(0.5, 0.3, f'(Based on avg: {monthly.mean():,.0f} MAD/month)',
             ha='center', va='center', fontsize=12)
    ax4.axis('off')
    
    plt.suptitle('Burn Rate Dashboard', fontsize=18, weight='bold', y=0.98)
    plt.tight_layout()
    plt.show()
\end{lstlisting}

\section{Plot 10: Year-over-Year Comparison}

\subsection{Purpose}
Compare same month across different years to identify seasonal patterns.

\subsection{Implementation}

\begin{lstlisting}[style=pythonstyle]
def plot_yoy_comparison(self):
    """Year-over-year comparison by month."""
    
    # Extract year and month
    self.df['year'] = self.df['month'].dt.year
    self.df['month_num'] = self.df['month'].dt.month
    
    # Aggregate: sum per year-month
    pivot = self.df.groupby(['year', 'month_num'])['salaire_brut'].sum().unstack(level=0)
    
    # Plot
    fig, ax = plt.subplots(figsize=(14, 7))
    
    colors = ['#2E86AB', '#A23B72', '#F18F01', '#06A77D']
    for i, year in enumerate(pivot.columns):
        ax.plot(pivot.index, pivot[year],
                marker='o', label=f'Year {year}',
                linewidth=2, color=colors[i % len(colors)])
    
    ax.set_title('Year-over-Year Payroll Comparison',
                 fontsize=16, weight='bold')
    ax.set_xlabel('Month', fontsize=12)
    ax.set_ylabel('Total Cost (MAD)', fontsize=12)
    ax.set_xticks(range(1, 13))
    ax.set_xticklabels(['Jan', 'Feb', 'Mar', 'Apr', 'May', 'Jun',
                        'Jul', 'Aug', 'Sep', 'Oct', 'Nov', 'Dec'])
    ax.legend()
    ax.grid(True, alpha=0.3)
    
    plt.tight_layout()
    plt.show()
\end{lstlisting}

\subsection{Insights}

\begin{itemize}
    \item \textbf{Seasonal spikes:} Bonuses in December
    \item \textbf{Year-over-year growth:} Hiring expansion
    \item \textbf{Consistent pattern:} Predictable costs
\end{itemize}
